\chapter{Backend}

\section{Estructura del Proyecto}

El backend sigue la estructura est\'andar de Maven con Spring Boot:

\begin{verbatim}
src/main/java/com/veterinary/
  |-- VeterinaryApplication.java       # Entry point @SpringBootApplication
  |-- config/                          # Configuraciones (4)
  |   |-- WebConfig.java               # Configuracion web y CORS
  |   |-- WebSocketConfig.java         # STOMP over SockJS
  |   |-- MailConfig.java              # JavaMailSender
  |   |-- JacksonConfig.java           # ObjectMapper personalizado
  |-- controller/                      # REST Controllers (20)
  |-- domain/                          # Entidades JPA (20 + 9 enums)
  |-- dto/                             # Data Transfer Objects (6)
  |-- repository/                      # Spring Data JPA (19)
  |-- scheduler/                       # Tareas programadas (1)
  |   |-- VacunaScheduler.java         # Alertas de vacunacion
  |-- security/                        # Seguridad JWT (4)
  |   |-- JwtTokenProvider.java        # Generacion/validacion JWT
  |   |-- JwtAuthFilter.java           # Filtro de autenticacion
  |   |-- SecurityConfig.java          # Reglas de acceso
  |   |-- CustomUserDetailsService.java # Carga de usuarios
  |-- service/                         # Servicios (22)
      |-- CitaService.java             # Gestion de citas
      |-- FacturaService.java          # Facturacion y pagos
      |-- PdfService.java              # Generacion de PDFs iText 7
      |-- ReporteService.java          # Reportes y exportacion Excel
      |-- NotificacionService.java     # Emails con plantillas HTML
      |-- PagoOnlineService.java       # Integracion Stripe
      |-- AuditService.java            # Registro de auditoria
      |-- ... (15 servicios mas)
\end{verbatim}

\section{Controladores REST}

El sistema expone 20 controladores REST que cubren todos los m\'odulos:

\subsection{Autenticaci\'on}

\texttt{AuthController} -- \texttt{/api/auth}
\begin{itemize}
    \item \texttt{POST /login}: Autentica personal o cliente, retorna JWT con claims (email, rol, userId).
\end{itemize}

\subsection{Gesti\'on Cl\'inica}

Todos los siguientes endpoints est\'an bajo \texttt{/api/medico/} y requieren autenticaci\'on JWT:

\texttt{DashboardController}
\begin{itemize}
    \item \texttt{GET /dashboard}: Estad\'isticas del sistema (clientes, mascotas, citas hoy, ingresos del mes, stock bajo).
\end{itemize}

\texttt{VeterinarioController}
\begin{itemize}
    \item CRUD completo de veterinarios (s\'olo ADMIN).
    \item \texttt{PUT /perfil}: Actualizar perfil propio.
    \item \texttt{POST /cambiar-password}: Cambiar contrase\~na con verificaci\'on de contrase\~na actual.
\end{itemize}

\texttt{ClienteController}
\begin{itemize}
    \item CRUD de clientes con b\'usqueda de texto libre (\texttt{?q=texto}) sobre nombre, apellido y email.
\end{itemize}

\texttt{MascotaController}
\begin{itemize}
    \item CRUD de mascotas con listado por cliente.
    \item Vista detallada con citas, historial y vacunas.
\end{itemize}

\texttt{CitaController}
\begin{itemize}
    \item CRUD de citas con validaci\'on de solapamiento.
    \item WebSocket broadcast en \texttt{/topic/citas} al crear/actualizar/eliminar.
    \item Notificaci\'on por email al crear la cita.
\end{itemize}

\texttt{ServicioController} -- Cat\'alogo de servicios con activaci\'on/desactivaci\'on.
\texttt{MedicamentoController} -- Inventario con ajuste de stock y alertas de stock bajo.
\texttt{VacunaController} -- Registro de vacunas por mascota con control de pr\'oxima dosis.
\texttt{HistorialMedicoController} -- Notas de consulta vinculadas a citas.
\texttt{HospitalizacionController} -- Ingreso/alta de mascotas con control de jaula.
\texttt{PlanTratamientoController} -- Planes con pasos secuenciados y estados.
\texttt{ConsentimientoController} -- Consentimientos con firma digital y generaci\'on de PDF.
\texttt{OdontogramaController} -- Diagrama dental con 42 dientes y 8 estados.
\texttt{FacturaController} -- Facturaci\'on con PDF, pagos parciales y control de saldo.
\texttt{ReporteController} -- Reporte de ingresos por per\'iodo con exportaci\'on a Excel.

\subsection{Portal Cliente}

\texttt{PortalClienteController} -- \texttt{/api/portal}
\begin{itemize}
    \item \texttt{GET /perfil}: Perfil del cliente autenticado.
    \item \texttt{GET /citas}: Citas del cliente.
    \item \texttt{GET /facturas}: Facturas del cliente.
    \item \texttt{GET /mascotas}: Mascotas del cliente con historial.
\end{itemize}

\subsection{P\'ublico}

\texttt{PublicBookingController} -- \texttt{/api/public} (sin autenticaci\'on)
\begin{itemize}
    \item \texttt{GET /veterinarios}: Lista de veterinarios activos con horarios.
    \item \texttt{GET /slots?vetId=\&fecha=}: Slots disponibles calculados por horario y citas existentes.
    \item \texttt{POST /citas}: Crear cita como visitante (validaci\'on de email contra BD).
\end{itemize}

\subsection{Pagos Online}

\texttt{PagoOnlineController} -- \texttt{/api/pagos}
\begin{itemize}
    \item \texttt{POST /stripe/create-payment-intent}: Crea intenci\'on de pago Stripe vinculada a una factura.
    \item \texttt{POST /stripe/webhook}: Webhook de confirmaci\'on de pago.
\end{itemize}

\subsection{Auditor\'ia}

\texttt{AuditController} -- \texttt{/api/admin/auditoria}
\begin{itemize}
    \item \texttt{GET /}: Listado paginado de logs de auditor\'ia ordenados por fecha descendente.
\end{itemize}

\section{Capa de Servicios}

22 servicios que encapsulan la l\'ogica de negocio. Los m\'as relevantes:

\begin{itemize}
    \item \texttt{CitaService}: Validaci\'on de disponibilidad (evita solapamientos usando JPQL), notificaciones por email y broadcast WebSocket.
    \item \texttt{FacturaService}: Creaci\'on autom\'atica de n\'umero de factura (FAC-yyyyMMdd-HHmmss-XXXX), descuento de stock al incluir medicamentos, control de pagos parciales.
    \item \texttt{MedicamentoService}: Ajuste de stock con registro autom\'atico de movimiento de inventario (trazabilidad completa).
    \item \texttt{NotificacionService}: Plantillas HTML para emails de recordatorios, alertas de vacunas y notificaciones de facturas.
    \item \texttt{PdfService}: Generaci\'on de PDFs con iText 7 (facturas con detalle de items, consentimientos informados).
    \item \texttt{ReporteService}: Reporte de ingresos con Chart.js y exportaci\'on a Excel con Apache POI.
    \item \texttt{PagoOnlineService}: Integraci\'on con Stripe Payment Intents en modo sandbox cuando no hay API key configurada.
    \item \texttt{PublicBookingService}: C\'alculo de slots disponibles basado en horarios de veterinarios y citas existentes.
    \item \texttt{AuditService}: Registro centralizado de operaciones CRUD con metadatos (usuario, IP, fecha).
\end{itemize}

\subsection{Ejemplo: Validaci\'on de Solapamiento de Citas}

\begin{lstlisting}[language=Java, caption=M\'etodo JPQL para detectar solapamiento]
@Query("SELECT c FROM Cita c WHERE c.veterinario.id = :vetId "
     + "AND c.estado NOT IN ('CANCELADA', 'NO_ASISTIO') "
     + "AND c.fechaHoraInicio < :fin AND c.fechaHoraFin > :inicio "
     + "AND (:citaExcluirId IS NULL OR c.id <> :citaExcluirId)")
List<Cita> findCitasOcupadas(@Param("vetId") Long vetId,
                             @Param("inicio") LocalDateTime inicio,
                             @Param("fin") LocalDateTime fin,
                             @Param("citaExcluirId") Long citaExcluirId);
\end{lstlisting}

\subsection{Ejemplo: Generaci\'on de PDF con iText 7}

\begin{lstlisting}[language=Java, caption=Generaci\'on de PDF de factura]
public ByteArrayInputStream generarFacturaPdf(Factura factura) {
    ByteArrayOutputStream out = new ByteArrayOutputStream();
    PdfWriter writer = new PdfWriter(out);
    PdfDocument pdf = new PdfDocument(writer);
    Document document = new Document(pdf, PageSize.A4);

    document.add(new Paragraph("PETCLINIC VETERINARY CLINIC")
        .setFontSize(20).setBold().setTextAlignment(TextAlignment.CENTER));
    document.add(new Paragraph("Factura: " + factura.getNumeroFactura())
        .setFontSize(14).setTextAlignment(TextAlignment.CENTER));
    document.add(new Paragraph(" "));

    document.add(new Paragraph("Cliente: "
        + factura.getCliente().getNombre() + " "
        + factura.getCliente().getApellido()));
    document.add(new Paragraph("Fecha: "
        + factura.getFechaEmision().format(
            DateTimeFormatter.ofPattern("dd/MM/yyyy HH:mm"))));

    Table table = new Table(new float[]{3, 1, 1, 1});
    table.addHeaderCell("Descripcion");
    table.addHeaderCell("Cant.");
    table.addHeaderCell("P. Unit.");
    table.addHeaderCell("Subtotal");
    for (DetalleFactura d : factura.getDetalles()) {
        table.addCell(d.getDescripcionLinea());
        table.addCell(String.valueOf(d.getCantidad()));
        table.addCell("$" + d.getPrecioUnitario());
        table.addCell("$" + d.getSubtotal());
    }
    document.add(table);
    document.add(new Paragraph("Total: $" + factura.getTotal())
        .setFontSize(16).setBold().setTextAlignment(TextAlignment.RIGHT));
    document.close();
    return new ByteArrayInputStream(out.toByteArray());
}
\end{lstlisting}

\subsection{Ejemplo: Integraci\'on con Stripe}

\begin{lstlisting}[language=Java, caption=Creaci\'on de Payment Intent en PagoOnlineService]
public String createPaymentIntent(Long facturaId) {
    Factura factura = facturaRepo.findById(facturaId)
        .orElseThrow(() -> new RuntimeException("Factura no encontrada"));

    // Stripe maneja montos en centavos
    long montoCentavos = (long)(factura.getTotal() * 100);

    PaymentIntentCreateParams params = PaymentIntentCreateParams.builder()
        .setAmount(montoCentavos)
        .setCurrency("mxn")
        .putMetadata("facturaId", facturaId.toString())
        .putMetadata("clienteId", factura.getCliente().getId().toString())
        .build();

    try {
        PaymentIntent intent = PaymentIntent.create(params);
        factura.setStripePaymentIntentId(intent.getId());
        facturaRepo.save(factura);
        return intent.getClientSecret();
    } catch (StripeException e) {
        throw new RuntimeException("Error al crear payment intent: "
            + e.getMessage());
    }
}
\end{lstlisting}

\section{Configuraci\'on}

\subsection{Application Properties}

\begin{lstlisting}[caption=Configuraci\'on principal]
server.port=8090

# Base de Datos PostgreSQL
spring.datasource.url=jdbc:postgresql://localhost:5432/veterinary_db
spring.datasource.username=veterinary_user
spring.datasource.password=veterinary_pass

# JPA / Hibernate
spring.jpa.hibernate.ddl-auto=update
spring.jpa.properties.hibernate.dialect=org.hibernate.dialect.PostgreSQLDialect
spring.jpa.show-sql=false

# JWT
app.jwt.secret=${JWT_SECRET:ClaveSuperSeguraParaFirmarTokensJWT2024}
app.jwt.expiration=${JWT_EXPIRATION:1800000}

# Stripe
stripe.api-key=${STRIPE_API_KEY:sk_test_placeholder}
stripe.webhook-secret=${STRIPE_WEBHOOK_SECRET:whsec_placeholder}

# Mail (SMTP)
spring.mail.host=${MAIL_HOST:smtp.gmail.com}
spring.mail.port=${MAIL_PORT:587}
spring.mail.username=${MAIL_USERNAME:}
spring.mail.password=${MAIL_PASSWORD:}
spring.mail.properties.mail.smtp.auth=true
spring.mail.properties.mail.smtp.starttls.enable=true

# Clinica
app.clinic.nombre=PetClinic Veterinary Clinic
app.clinic.email=contacto@petclinic.com
\end{lstlisting}

\subsection{DataInitializer}

Al arrancar por primera vez, \texttt{DataInitializer} inserta datos de prueba utilizando el patr\'on \texttt{CommandLineRunner}:

\begin{itemize}
    \item 3 veterinarios: Admin Sistema (ADMIN, admin@petclinic.com/Admin123), Laura Garc\'ia (VETERINARIO, Cirug\'ia, vet@petclinic.com/Vet123), Carlos L\'opez (RECEPCIONISTA, recepcion@petclinic.com/Recep123).
    \item 2 clientes: Ana Mart\'inez (ana@email.com, con portal activo/Ana123), Pedro Ram\'irez (pedro@email.com).
    \item 3 mascotas: Max (Golden Retriever, 30kg), Luna (Siam\'es, 4kg), Rocky (Bulldog Franc\'es, 12kg).
    \item 5 servicios: Consulta General (\$500), Vacunaci\'on (\$350), Cirug\'ia (\$2,500), Desparasitaci\'on (\$200), Grooming (\$400).
    \item 3 medicamentos: Amoxicilina 500mg (100 tabs, stock m\'inimo 10), Meloxicam 5mg (50 ml, stock m\'inimo 5), Frontline Plus (3 pipetas, stock m\'inimo 2).
    \item 2 citas programadas para el d\'ia siguiente con estados PROGRAMADA y CONFIRMADA.
\end{itemize}

\section{Manejo de Excepciones}

El \texttt{GlobalExceptionHandler} captura y formatea errores de manera consistente retornando JSON con estructura \texttt{\{error: mensaje, details: [...]\}}:

\begin{longtable}{|p{5cm}|p{4cm}|p{5cm}|}
\hline
\textbf{Excepci\'on} & \textbf{C\'odigo HTTP} & \textbf{Uso} \\
\hline
EntityNotFoundException & 404 & Recurso no encontrado (ej: cliente, mascota, factura) \\
\hline
IllegalArgumentException & 400 & Par\'ametros inv\'alidos (ej: fechas incorrectas, datos duplicados) \\
\hline
AccessDeniedException & 403 & Sin permisos para el recurso solicitado \\
\hline
MethodArgumentNotValidException & 400 & Errores de validaci\'on de DTOs (\texttt{@NotBlank}, \texttt{@Email}, etc.) \\
\hline
RuntimeException & Variable & Determinado por el mensaje (ej: ``Stock insuficiente'' $\rightarrow$ 400) \\
\hline
Exception & 500 & Error interno del servidor \\
\hline
\end{longtable}
